\subsection{Prediction 8: Why Good Incentives Go Bad}
\label{subsec:pred8-mechanism}

%%%%%%%%%%%%%%%%%%%%%%%%%%%%%%%%%%%%%%%%%%%%%%%%%%%%%%%%%%%%%%%%%%%%%%%%%%%%%%%
% CHAPTER 11.8: Narrative Treatment of Prediction 8
% Mechanism Design Must Model Context Dynamics
%%%%%%%%%%%%%%%%%%%%%%%%%%%%%%%%%%%%%%%%%%%%%%%%%%%%%%%%%%%%%%%%%%%%%%%%%%%%%%%

\subsubsection{The Day Care Disaster}

In 1998, economists Uri Gneezy and Aldo Rustichini conducted a famous \
experiment in Israeli day care centers.

The problem was simple: parents were picking up their children late, \
forcing teachers to stay overtime. The economists proposed a simple \
solution: a fine. Charge parents \$3 for every late pickup.

Standard economics predicted the outcome: higher cost → less lateness.

What actually happened: lateness \emph{doubled}.

The fine transformed a social relationship (``I'm inconveniencing people \
who care for my child'') into a market transaction (``I'm paying for \
extra service''). Parents felt \emph{less} guilty after paying. The moral \
constraint dissolved into a mere fee.

And here's the kicker: when the researchers removed the fine, lateness \
stayed high. The original social norm was gone, and wouldn't come back.

\subsubsection{The Mechanism Design Blind Spot}

Standard mechanism design assumes a fixed utility function: $U(\theta)$, \
where $\theta$ represents the agent's ``type.'' The designer's job is \
to create incentive-compatible mechanisms---structures where revealing \
true preferences is optimal.

But the day care experiment reveals something the standard theory misses: \
\emph{the mechanism changes the utility function it's designed for}.

Before the fine:
\begin{equation}
U = U_{practical} + U_{social} + U_{moral}
\end{equation}

After the fine:
\begin{equation}
U = U_{practical} + U_{market}
\end{equation}

The social and moral components didn't just shrink---they transformed \
into something else entirely.

\subsubsection{The Framework's Prediction}

Our framework captures this with the context-dependent utility function:

\begin{equation}
U = U(\theta, M, \Psi(M))
\end{equation}

where $M$ is the mechanism and $\Psi(M)$ is the context induced by the \
mechanism.

The critical insight: $\frac{\partial \Psi}{\partial M} \neq 0$.

Mechanisms don't just operate in contexts---they create them. And the \
context they create may violate the assumptions they were built on.

\subsubsection{How Mechanisms Change Context}

Consider the Ψ-dimensions affected by introducing incentives:

\paragraph{$\Psi_S$ (Social Context).}

Financial incentives signal that the interaction is market-like, not \
relationship-like. The blood donation example is canonical: paying people \
to donate blood \emph{reduces} donations, because it transforms a gift \
into a sale.

\paragraph{$\Psi_M$ (Moral Framing).}

Incentives can crowd out moral motivation. If I'm paid to do something \
good, it's no longer as good---the act becomes instrumental rather than \
expressive.

\paragraph{$\Psi_I$ (Institutional Trust).}

Introducing monitoring and incentives signals distrust. ``If they're \
watching me, they must think I'd cheat.'' This can become self-fulfilling.

\paragraph{$\Psi_T$ (Time Orientation).}

Short-term incentives focus attention on immediate payoffs, potentially \
crowding out long-term considerations that intrinsic motivation would \
have supported.

\subsubsection{Three Case Studies}

\paragraph{Case 1: Teacher Performance Pay.}

The theory: Pay teachers based on student test scores. Better performance \
→ higher pay → better outcomes.

The reality: Teachers ``teach to the test,'' narrow the curriculum, \
sometimes cheat. Intrinsic motivation (``I teach because I care about \
kids'') gets crowded out by extrinsic motivation (``I teach what's \
measured'').

The mechanism changed $\Psi_M$ (from calling to job) and $\Psi_T$ (from \
long-term development to short-term scores).

\paragraph{Case 2: Carbon Offsets.}

The theory: Create a market for emissions reductions. Efficient price \
discovery → optimal abatement.

The reality: Offsets often become permission slips rather than \
complements to reduction. Flying becomes ``guilt-free'' because you \
paid \$20. Total emissions don't fall.

The mechanism changed $\Psi_S$ (from collective responsibility to \
individual transaction) and reduced the coherence cost of high-emission \
behavior.

\paragraph{Case 3: Gig Economy Ratings.}

The theory: Customer ratings discipline service quality. Bad drivers \
get deactivated → good service.

The reality: Drivers game the system, offer bribes for good ratings, \
refuse difficult rides. The adversarial relationship between platform \
and workers undermines the cooperation the system needs.

The mechanism changed $\Psi_I$ (from employment relationship to \
surveillance relationship).

\subsubsection{The Γ-Matrix Implication}

The framework quantifies these effects through the full $8 \times 6$ \
Γ-matrix. Mechanism designers must account for:

\begin{center}
\renewcommand{\arraystretch}{1.3}
\begin{tabular}{ll}
\hline
\textbf{Mechanism Choice} & \textbf{Context Effect} \\
\hline
Financial incentives & $\Delta\Psi_M < 0$, $\Delta\Psi_S < 0$ \\
Monitoring & $\Delta\Psi_I < 0$ \\
Rankings/competition & $\Delta\Psi_S$ ambiguous \\
Public commitment & $\Delta\Psi_M > 0$, $\Delta\Psi_I > 0$ \\
\hline
\end{tabular}
\end{center}

The key parameter is $\frac{\partial \Psi}{\partial M}$---mechanism-context \
sensitivity. When this is high, standard mechanism design fails.

\subsubsection{Testable Implication}

\begin{tcolorbox}[colback=blue!5!white, colframe=blue!75!black, title=Prediction 8: Mechanism-Context Dynamics]
Incentive schemes that ignore context dynamics will systematically \
underperform their theoretical predictions, and may produce perverse \
effects that persist even after the incentives are removed.
\end{tcolorbox}

Specific test: The ``motivation crowding'' effect will be strongest in \
domains with high baseline $\Psi_M$ (moral framing) and high baseline \
$\Psi_S$ (social embeddedness). Financial incentives will backfire most \
spectacularly where intrinsic motivation was strongest.

\subsubsection{Design Implications}

The framework suggests principles for context-aware mechanism design:

\begin{enumerate}
    \item \textbf{Measure baseline Ψ.} Before introducing incentives, \
    understand the existing context. High $\Psi_M$ or $\Psi_S$ signals \
    crowding-out risk.
    
    \item \textbf{Use incentive types strategically.} Non-monetary \
    recognition may preserve $\Psi_M$ better than cash. Public commitment \
    may strengthen rather than weaken social motivation.
    
    \item \textbf{Test for irreversibility.} The day care study shows that \
    context changes can be permanent. Pilot programs should check whether \
    removing the mechanism restores the original context.
    
    \item \textbf{Model the full Γ-matrix.} Don't just ask ``will this \
    incentive work?'' Ask ``how will this incentive change the utility \
    function it operates on?''
\end{enumerate}

\subsubsection{The Deeper Point}

Samuel Bowles titled his book on this topic ``The Moral Economy.'' His \
argument: economic policy has moral side effects that economists \
systematically ignore.

Our framework formalizes this insight. The Ψ-dimensions are \emph{policy \
variables}, not just background conditions. Every mechanism is also an \
intervention in the moral and social fabric of the context it operates in.

This is prediction 8: mechanism design that ignores context dynamics \
isn't just incomplete---it's potentially self-defeating.

The day care fine didn't just fail to reduce lateness. It destroyed the \
norm that would have reduced lateness without any fine at all.

\begin{cross-reference}
\textbf{Formal Specification:} PRED-08 with mathematical treatment in \
Appendix AU, Section~\ref{subsec:pred-08-calc} \
(Novel Predictions: Formal Specification).
\end{cross-reference}
